\par\medskip
\subsection{Respuesta térmica durante la ocupación de oficinas}
La temperatura operativa se compara en las seis oficinas intervenidas mediante 3.130 intervalos horarios ocupados por zona: lunes a sábado, de 08:00 a 18:00, según el calendario 2025 del modelo, sin festivos adicionales. Se conservan los mismos horarios en ambas corridas. El percentil 95 describe la parte alta de la distribución; no es un umbral normativo de confort.
\begin{table}[H]\centering
\caption{Temperatura operativa en horas ocupadas: caso base y C3 revisión 02}
{\small\singlespacing\renewcommand{\arraystretch}{1.25}
\begin{tabular}{lrrrrr}\toprule
 & \multicolumn{2}{c}{Media (\textdegree C)} & \multicolumn{2}{c}{P95 (\textdegree C)} & \\
Oficina & CB & C3 r02 & CB & C3 r02 & $\Delta$ media (K) \\ \midrule
PB-OF-01 & 26,75 & 26,02 & 28,36 & 27,04 & -0,74 \\
PB-OF-02 & 26,98 & 26,14 & 28,56 & 27,24 & -0,85 \\
PB-OF-03 & 26,83 & 26,08 & 28,34 & 27,10 & -0,75 \\
1PA-OF-01 & 27,17 & 26,28 & 28,91 & 27,68 & -0,89 \\
1PA-OF-02 & 27,07 & 26,26 & 28,15 & 27,24 & -0,80 \\
1PA-OF-03 & 26,83 & 25,97 & 28,23 & 27,21 & -0,85 \\
\bottomrule\end{tabular}}
\fuente{Elaboración propia con series horarias de EnergyPlus y horarios de ocupación del modelo. $\Delta$ = C3 r02 $-$ CB.}
\end{table}
Estos estadísticos cuantifican la respuesta térmica, pero no acreditan C01. La evaluación de confort requiere separar las horas de funcionamiento mecánico y natural, aplicar el método correspondiente y comprobar sus condiciones de actividad, vestimenta, velocidad del aire y cobertura espacial. Una reducción de demanda o de temperatura media no equivale automáticamente a mayor porcentaje de horas confortables.
